\section{Introduction}
\label{sec:introduction}

Consider a household robot waiting for its next request.
Before a user asks for anything, it may inspect rooms, verify object locations, or acquire views that reveal attributes and spatial relations.
The identity of the future request is not known: it may ask the robot to navigate to a mug, report the color of a chair, determine whether a medication box is inside a cabinet, or verify that no extinguisher is present in a region.
Nor is the waiting period fixed.
A request can arrive after ten seconds or several minutes and can interrupt the robot halfway through a long traversal.
The robot must then serve the request using exactly the pose and evidence committed at that moment.

This regime is poorly captured by the dominant \emph{explore-then-act} abstraction.
Task-agnostic exploration typically optimizes geometric coverage, map completeness, information gain, or a checkpoint score after a fixed budget
\citep{yamauchi1997frontier,chaplot2020ans,ye2026look,zhang2026theory}.
Pre-task context and long-term memory methods make the acquired state reusable
\citep{cai2026autocontext,wang2026lmee,jiang2024roboexp,chen2026abot},
but they likewise evaluate a completed exploration phase or the quality of its terminal memory.
Task-conditioned active perception instead values information through a known goal
\citep{liu2025activevoo,lee2026obsgraph}.
Neither view answers the decision faced while the task is still hidden and the exploration action may be interrupted.

The missing variable is \emph{when evidence becomes usable}.
Suppose two actions expose the same facts after 100 low-level steps.
The first action reveals a frequently requested object after 20 steps and then continues to a distant region; the second reveals nothing until its final panoramic scan.
They are equivalent to any endpoint-only objective.
They are not equivalent if most requests arrive before step 50.
Macro-actions such as \emph{go to the kitchen and scan the counters} are not atomic, so endpoint scoring can systematically prefer actions with high terminal value but poor early prefixes.

We formalize this problem as \textbf{interruptible pre-task exploration} (\taskname).
Before exploration, the agent knows a distribution over a catalog of possible tasks and an arrival hazard, but not the realized task or arrival time.
At every completed low-level step, its sensor observation is atomically committed to memory and an arrival event is checked.
If a request arrives, exploration halts immediately.
Performance is the true downstream utility of the frozen memory and pose, optionally after a small post-arrival service budget.
This formulation evaluates \emph{anytime readiness} rather than terminal exploration quality.

We propose \method{}---the \textbf{C}ounterfactual \textbf{A}nytime \textbf{R}eadiness \textbf{A}gent.
Its central component is a \textbf{joint evidence--time model} (\jet{}).
For each candidate macro-action, \jet{} predicts a joint distribution over the first prefix at which reusable evidence primitives become reliable.
The primitives are stored in a reliability-first memory (\rfm{}) with provenance, view quality, timestamps, and calibrated confidence.
A task template is compiled into a requirement program over these primitives.
Executing each program on joint evidence-time samples produces a task readiness curve without conditioning exploration on the eventual task.
An arrival-aware semi-Markov planner integrates these curves against the request hazard and trades interruption utility against action cost and future value.
Counterfactual simulator branches provide supervision at real action prefixes, including right-censored actions that never acquire a particular primitive.

\begin{figure}[t]
  \centering
  \placeholderfigure{3.8cm}{Introduction overview: fixed-budget exploration versus \taskname{}; a request interrupts a macro-action; \method{} prioritizes evidence that becomes useful early.}
  \caption{\textbf{Ready before asked.}
  Conventional exploration scores the endpoint of a fixed budget.
  \taskname{} evaluates the downstream service supported at a random interruption time.
  \method{} predicts evidence acquisition times and selects actions that maximize arrival-weighted readiness.}
  \label{fig:overview}
\end{figure}

We instantiate the setting in \bench{}.
A controlled Mini-\taskname{} layer exactly enumerates tasks, arrivals, prefixes, and optimal policies, isolating the mechanism from perception.
The embodied evaluation then spans procedurally generated ProcTHOR houses, scanned HM3D environments, high-quality HSSD interiors, Habitat-GS Gaussian-splatting scenes, and dynamic HSSD histories from P4D-Bench
\citep{savva2019habitat,ramakrishnan2021hm3d,deitke2022procthor,khanna2024hssd,xia2026habitatgs}.
HM3D-OVON provides open-vocabulary navigation, OpenEQA supplies natural attribute and spatial questions, HSSD ObjectNav tests synthetic-to-scanned transfer, Habitat-GS PointNav isolates renderer and control transfer, and P4D-Bench stresses evidence that becomes stale across sessions
\citep{yokoyama2024hm3dovon,majumdar2024openeqa}.
The native task definitions are retained where applicable, while the realized task is withheld until a sampled interruption time.
All methods share RGB-D sensing, noisy odometry, frozen perception and navigation interfaces, memory capacity, and downstream executors within each domain.
Four task families---category/instance navigation, attribute QA, spatial QA, and presence/absence verification---share an explicit evidence vocabulary and include held-out compositions.
The evaluation spans five arrival processes and real indoor/outdoor deployment on DEEP Robotics LYNX M20, X30, and Lite3 robots with a shared high-level interface.

Our contributions are threefold:
\begin{enumerate}[leftmargin=*,itemsep=2pt,topsep=3pt]
  \item \textbf{Problem and protocol.}
  We introduce \taskname{}, in which a hidden task can arrive inside a macro-action, and measure the downstream utility supported by the exact interrupted prefix.
  \item \textbf{Method.}
  We introduce \method{}, which learns joint first-hitting times of reliable evidence from counterfactual action prefixes, composes them with task requirement programs, and plans with an arrival-aware semi-Markov objective.
  \item \textbf{Benchmark and diagnosis.}
  We define \bench{} across controlled, procedural, scanned, synthetic, Gaussian-splatting, dynamic, and real-world environments and demonstrate physical deployment on three robots, testing endpoint-proxy mismatch, timing sensitivity, cross-domain transfer, compositional generalization, and robustness to arrival and task-prior shift.
\end{enumerate}
